\chapter{Depression}

\section*{Definition}
Depression is a condition involving persistent low, sad, or irritable mood and/or loss of interest or pleasure, together with other symptoms that cause distress or impair functioning. A depressive episode commonly lasts at least two weeks, but assessment considers the whole clinical course.

\section*{When to See a Doctor}
See your doctor if:
\begin{itemize}
  \item Mood symptoms persist beyond two weeks or you have thoughts of self-harm
\end{itemize}

\section*{How It Develops}
Depression involves interacting biological, psychological, and social factors; it is not adequately explained by a simple “chemical imbalance.” Assessment distinguishes it from bipolar disorder, grief, adjustment reactions, substance effects, and medical conditions.

\section*{Causes}
\begin{itemize}
  \item Chemical imbalance in brain
  \item Life events trauma genetic predisposition
  \item Chronic medical conditions
  \item Medication side effects
  \item Substance use disorders
\end{itemize}

\section*{Related Symptoms}
\begin{itemize}
  \item Sleep disturbance
  \item Appetite weight changes
  \item Fatigue/loss of energy
  \item Feelings worthlessness guilt
  \item Thoughts death suicide
\end{itemize}


\section*{Medical Evaluation}
\begin{itemize}
  \item Comprehensive psychiatric assessment includes history of present illness, past psychiatric history, substance use, medical history, and collateral information from family.
  \item A validated measure such as the PHQ-9 may support assessment and follow-up but does not replace a diagnostic interview; tools for anxiety or trauma are used when relevant.
  \item Targeted medical tests are guided by symptoms, examination, medicines, and history rather than ordered routinely for everyone.
  \item Risk assessment for self-harm, suicide, and danger to others is an essential component of every psychiatric evaluation.
\end{itemize}

\section*{Treatment Options}
\begin{itemize}
  \item Evidence-based psychotherapy, including cognitive-behavioral or interpersonal therapy, is a first-line treatment option for depressive disorders; the choice depends on severity, preference, risk, and comorbidity.
  \item Pharmacotherapy with antidepressants, anxiolytics, mood stabilizers, or antipsychotics is tailored to the specific diagnosis and individual patient factors.
  \item Combined treatment may be appropriate for some people, especially with more severe or recurrent depression; treatment should reflect severity, preference, risk, and comorbidity.
  \item Referral to psychiatry is indicated for severe symptoms, treatment resistance, suicidal ideation, or diagnostic uncertainty.
\end{itemize}

\section*{Self-Care and Home Management}
\begin{itemize}
  \item Establish a consistent daily routine with regular sleep, meals, and physical activity to support mental health stability.
  \item Practice stress management techniques including mindfulness, deep breathing exercises, and progressive muscle relaxation.
  \item Maintain social connections and avoid isolation by reaching out to trusted friends, family, or support groups.
  \item Avoid alcohol, caffeine, and recreational drugs which can destabilize mood and exacerbate psychiatric symptoms.
  \item If experiencing suicidal thoughts or crisis, contact emergency services or a crisis helpline immediately.
\end{itemize}

\section*{References}
\begin{thebibliography}{99}
\bibitem{depression:ref1} American Psychiatric Association. ``Major depressive disorder.'' In: Diagnostic and Statistical Manual of Mental Disorders. 5th ed, Text Revision. APA; 2022.
\bibitem{depression:ref2} Malhi GS, Mann JJ. ``Depression.'' Lancet. 2018;392(10161):2299-2312.
\bibitem{depression:ref3} Gelenberg AJ, Freeman MP, Markowitz JC, et al. ``Practice guideline for the treatment of patients with major depressive disorder.'' 3rd ed. Am J Psychiatry. 2010;167(10 Suppl):1-152.
\bibitem{depression:ref4} Cipriani A, Furukawa TA, Salanti G, et al. ``Comparative efficacy and acceptability of 21 antidepressant drugs for the acute treatment of adults with major depressive disorder.'' Lancet. 2018;391(10128):1357-1366.
\bibitem{depression:ref5} Kupfer DJ, Frank E, Phillips ML. ``Major depressive disorder: new clinical, neurobiological, and treatment perspectives.'' Lancet. 2012;379(9820):1045-1055.
\bibitem{depression:ref6} Fava M. ``Unipolar depression in adults: assessment and diagnosis.'' UpToDate. Wolters Kluwer; 2025.
\end{thebibliography}
